\documentclass[11pt]{article}
\title{Response Memo}
\begin{document}
\maketitle

I thank all of the reviewers for what proved to be a very helpful, if challenging, set of comments and criticisms.  I believe that, thanks to this feedback, the present paper is much improved over the previous version.  In  what follows I respond to each of the major criticisms that I distilled from the reviewer reports, outlining the changes that I made to paper in response to those concerns where applicable.

\section{Reviewer 1}
\begin{enumerate}
	\item \emph{ {\ldots}a limitation of the data (\textbf{note-there is not a section on limitations in the manuscript}) is knowing the strength and frequency of interaction between employer and employee. Perhaps they know each other in some capacity. The analysis assumes they are distinct and that there is no relation among them.}
	\item \emph{A second concern is the assumption that popular cultural tastes are the content used by individuals in sustaining relationships with their "weak ties." While the conversational "tool-kit" model of cultural capital suggests that popular culture is used in conversations with weak ties, there is no evidence that actors actually did this in their contact with weak ties in efforts to find a job.  The model implies a moderate relationship, but does not demonstrate the actual use of cultural capital to manage these interpersonal relationships nor as a mechanism for activating instrumental relations.} 
	\item \emph{Finally, a limitation is that this data is a cross-sectional snap shot of respondents collected nine years ago. If we accept the argument in the paper, are things still operating the same way today? Repeating this analysis using more current data and discovering similar results would strengthen the theoretical arguments in the paper and offer something new to cultural sociologists. }
	
The reviewer gets at three important limitations of of the current effort. I agree with the reviewer that a section on limitations would do a lot to circumscribe the strengths and limits of the paper.  In this version of the paper I have included such a section, and in it I tackle some of the issues raised by the reviewer.   

In terms of the first two points the reviewer is of course quite correct in noting that there is no direct information about the nature of the employer/employee relationship, nor is there any direct information about the particular sort of cultural capital that was mobilized to activate the relationships in question.   This is unfortunately a limitation of relying on survey data---although all of the information at hand such as the higher likelihood of those who use acquaintances to find out about their current job to consume popular culture, points in that direction. The only way to really get at the issue of whether the specific mechanisms that I point to are operative on the ground is via qualitative and observational research. One can imagine constructing survey items for this purpose, but they are unlikely to fully satisfy the concerns pointed at by  the reviewer. I point to this issue during the conclusion.  

In terms of the last point, while I understand where the reviewer is coming from, I don't see any \emph{a priori} reason for why there should be any dramatic change in how the mechanisms that I outline here, which are fairly generic, operate in less than a decade.  Although I do wish---with the reviewer---that there was more recent data unfortunately the 2002 GSS was the last wave of this survey to include culture consumption items.
\end{enumerate}

\section{Reviewer 2}
\begin{enumerate}
	\item \emph{The paper's theory sections are too long. The paper makes a simple empirical advance that is easily situated with respect to earlier theoretical work by DiMaggio, Erickson, Peterson, Granovetter and Bourdieu. An involved theory section is not necessary.}
	
	I appreciate the reviewer's advice in this regard.  The current version of the paper sports a much more succinct theory section devoted exclusive to setting up the main argument under consideration.
	
	\item \emph{Section 2 on network ties as social capital tries to describe a distinction between two ways social ties are conceived as valuable, and I do not think this distinction really exists.} 
		
		The reviewer's point is well-taken.  I was using the distinction mostly as a way to organize the review of the social capital literature, not necessarily to make a conceptual point.  I agree with the reviewer that making this differentiation is orthogonal to the main argument of the paper.  In this version of the paper I review the social capital literature in a single section, and followed the reviewer's advice in characterizing both diversity and resources as being two components of a single process.
	
	\item \emph{The analysis presented on page 12 and in Table 2 is cumbersome. The author is trying to show that mean number of cultural forms participated in varies with how a person first found out about their current job. I think ANOVA would be the most straightforward test. How respondents learned about their current job divides the respondents into 7 categories. One-way ANOVA of number of cultural forms participated in across these categories will let you test whether this grouping of respondents accounts for a significant portion of the variance in number cultural forms participated in.  Equivalently, you can regress number of cultural goods consumed on 6 dummies for 6 of the 7 categories.}
	
	I agree with the reviewer that the previous way of presenting the results (first univariate then multivariate) was unnecessarily cumber and involved and did not match well with the way that I stated the hypotheses.  In this version of the paper I decided to jump straight to the multivariate analysis for the sake of succinctness.  So I eliminated this first data-analytic step.
	
	\item \emph{To demonstrate the first claim, the author needs to define a variable coded 1 if the respondent found her/his job through a tie. So if the respondent found the job through a relative, a friend, or an acquaintance or if the employer contacted the respondent directly, the variable is coded 1; otherwise 0.  The author does not do this. He just models the likelihood of using a particular type of tie relative to other types of tie and formal channels combined. For some types of tie (acquaintance and contacted by employer), the author finds that cultural variety positively affects use of this type of tie relative to everything else, but for the other two types of tie (relatives and close friends), he finds no significant effect. We cannot infer from these results that claim 1 is true.  To demonstrate the second claim, the author needs to conduct an analysis of just respondents who found their current job through a social tie and divide these respondents into a strong tie category (found job through a relative or friend) and weak tie category (found job through an acquaintance or was contacted by employer). Then regress found-job-through-weak-tie on cultural variety and controls.}
	
	I really appreciate the reviewer's point in this respect. The previous version of the paper did not display a tight match between the claims at stake and the procedures that I used to test them.  In this version of the paper I state the hypotheses that the reviewer is referring to explicitly (they are the main hypotheses---1 and 2---in this version of the paper).  I also followed the reviewer's obviously sound advice in terms of constructing the outcome variables.  Hypothesis one now pertains to the first outcome mentioned by the reviewer (network contact versus other) and the second outcome now pertains to--among those who used a contact to find out about their job--the weak versus strong tie comparison excluding those who did not use a social contact to find out about their job.   I believe that the reviewer will agree that this version of the paper now evinces a much tighter match between hypotheses (as stated) and the analytic procedures used to verify them.
	
	\item \emph{The author says that in earlier work using this same 2002 GSS data (Lizardo 2006), he found that network size and number of cultural activities participated in were positively related (page 8-9). So it seems like he can test for the link missing from all of this: Is network size positively related to having used a weak tie to get one's current job? If so, we can have much more confidence in the author's argument. If not, the author's interpretation of his results probably needs revision.}	
	
	Once again, I appreciate the reviewer's very good point.  This version of the paper now includes an analysis of the sort recommended by the reviewer and, as the reviewer will appreciate, the results broadly support the argument.
	
	\item \emph{It is strange to say that the respondent "activated" a tie to find out about a job when the respondent was directly contacted by the employer initially. The respondent did not activate that tie; the employer did. It still would be reasonable to say in this case that the respondent found out about her job (obtained her job) through a network tie.}
	
	I agree with this observation.  In this version of the paper I have rectified my terminology where appropriate.	
	
	\item \emph{On page 7, the author claims that network theories of social capital treat network ties as exogenous and provides no cites. I think there are influential papers that do not do this (DiMaggio 1987; Erickson 1996). There are also studies on gender and the life course that look at the determinants of ego network size and composition (McPherson and Smith-Lovin 1982, AJS; Moore 1990, ASR; Munch et al 1997, ASR).[\ldots]Maybe the author should be more precise and provide cites or avoid this claim. I don't see how this claim is critical to what the paper accomplishes.}
	
	The point that I wanted to convey is that a lot of network research does not deal with the factors that lead to the maintenance or formation of ties. I this version of the paper, I have reworded this claim accordingly.
	
	\item \emph{First sentence in section 4: Do you mean density or diversity? (Consumption of a variety of cultural goods is probably also related to density, but negatively.)}
	
	I meant diversity; I have changed the wording accordingly.	
\end{enumerate}

\section{Reviewer 3}
\begin{enumerate}
	\item \emph{{\ldots}from an explanatory point of view it is unclear the origin of that relationship. The author mentions the conversion of cultural capital into social capital, but conversions of capital can be both ways, and in Bourdieu's theory they are quite intermingled. So a theoretical challenge would be to disentangle that association. In order to do that, we would need a different methodological strategy that the one used, that is, it goes beyond the data analysis of a single cross-sectional survey. }
	
	The reviewer's point is of course precisely on the mark.  However, I think that the reviewer exaggerates the limits of cross-sectional data in this regard. Even when it comes to bi-directional processes, there are ways to get at the issue of causal direction.  In fact Lizardo (2006) attempts to do just that, and this research builds on that effort. 
	
	\item \emph{If we go back to the explanation, we could do the regressions but this time social capital would be the explanatory factor and cultural capital the explained one, and it would be interesting to look at the consistency of the results. }
	
	Not sure why the reviewer might think that these analyses would reveal anything dramatically different, but I assure the reviewer that the results are consistent even when cultural capital is the response and job-search method is the predictor.  I can make those results available on request.
	
	\item \emph{This first comment is important as it would require to work a bit more on the categories that are mentioned in the paper but never much developed statistically: popular culture, "high prestigious arts", and omnivorousness. Above all, this latter is mentioned but not statistically worked. There is a measure of volume, but letting it be as it is makes impossible to distinguish between the effects of these types offered. Instead of volume, it would be probably more accurate to distinguish between these three groups of people, and using these in the regressions run. This does not imply to turn this paper into one within the debate on the omnivore, which is already an area in itself. }
	
	I agree with the reviewer that entering into the omnivore debate would take this paper too far afield.  This is why I decided to concentrate on the more neutral term ``cultural variety'' which we owe to Erickson (1996).  I can see how splitting the activities by ``brow-level'' might reveal some important insights, but I decided against it because it would make the analysis too complicated.  I would be more open however if the reviewer were to specify a more theoretically driven reason for doing this.  In the absence of strong theory, I opted to go with the simplest way of characterizing the cultural capital variable which is the volume measure.  Note that I do disaggregate the activities when looking at differences between those who used acquaintances and those who were contacted directly by their employer and this analysis demonstrates that these are two very different types of culture consumers however.  
	
	\item \emph{I missed factors as trajectory or social mobility, as it would probably explained at least partially the effects of cultural capital. I understand that this might be beyond the questions of the survey analysed.}
	
The reviewer makes an excellent point.  The GSS does contain information on respondent's class background (e.g. parental education).   While parental education does not affect the overall probability of a respondent having found their current job via a network contact, it does have a strong positive effect on the probability that \emph{conditional} on having found your current job via a network contact, you used an acquaintance or were directly contacted by the employer rather than using a strong tie.  This is indicative of the fact that being directly contacted by the employer is an elite route to finding a job, but using a strong tie is generally correlated with markers of low status. The reviewer is right that including the parental education variable in the regression does partially ``explain'' the effect of cultural variety on finding a job via weak versus a strong tie, although I would be reticent to use the word ``explain'' here, since parental education precedes cultural capital (and cultural capital ``intervenes'' between parental education and social capital), so that it is more accurate to say that class background has both direct and indirect (mediated via its own positive effects on cultural variety) effects on social capital.  However, since in the paper I am not directly concerned with decomposing the effects of class background on social connectivity but instead on noting the direct effects of cultural capital on job search outcomes, I opted to not deal with this issue directly in the paper (the results are available on request).
\end{enumerate}

\section{Reviewer 4}
\begin{enumerate}
\item 	\emph{The writing might be improved in parts - it is very dry and abstract in places, and most of section 2.2 is incomprehensible. There are a fair number of typos too.}

I have tried to make this version of the paper more readable and succinct than the previous version.  I have also subjected it to a thorough proof-reading.

\item \emph{It wasn't clear that the final sentence of the paper which attributed the effects to face-to-face interaction was justified. I can't find reference to the survey question about contacts in the job search, but I would be surprised if the definition of 'knowing someone' specified that it was necessarily face-to-face. Could not an employer 'know' the person they contact by reputation rather than personal contact?}

The reviewer is of course correct in this regard.  Unfortunately, these data do not allow us to ascertain whether ``knowing'' refers to reputation or face-to-face from the survey.  However, secondary research on labor markets (see for instance Yakubovich 2005) suggests that \emph{the majority} of the cases of the employer ``knowing'' the prospective employee are in fact cases of direct personal knowledge.    So while it is true that the category ``knowing directly'' might include some cases where ``knowing'' does not refer to personal acquaintance it is unlikely the category is radically vitiated by cases of spurious knowing. 

\item \emph{I don't suppose it is possible to derive from the data, but it would be interesting to know what were typical jobs obtained through employer approach of mobilisation of acquaintances.}

The numbers do not allow for a full disaggregation of occupational categories, but as I note in the paper, on average the jobs obtained via direct employer contact or acquaintances tend to be (on average) more prestigious than jobs obtained via strong ties (close friends and relatives).  This is consistent with previous research on socio-economic status differences in labor market strategies.   The data do show however, that the most prestigious jobs are actually obtained via organizational mechanisms: employment services or schools.
\end{enumerate}
\end{document}•
